\newcommandx{\supD}[3][1=P, 2=\setQ, 3=X]{\sup_{Q \in #2} \kl[Q][#1][#3]}
\newcommand{\myDiv}[1][P]{\measure{D}{\mathrm{C}}{\setQ \,\middle\|\, #1(X)}}
\newcommandx{\klExp}[3][1=Q, 2=P, 3=x]{\sum_{#3} #1(#3) \ln \frac{#1(#3)}{#2(#3)}}
\newcommand{\inSet}[1]{\mathFix{#1 \in \set{#1}}}

\newcommand{\saddle}{\mathFix{(\saddP, \saddR)}}
\newcommand{\fExp}{\sum_i R(i) \klExp[Q_i]}
\newcommand{\cDist}[1][]{\mathFix{\sub{C}{#1}}}
\newcommand{\ccDist}[1][]{\mathFix{\sub{C^{*}}{#1}}}
\newcommand{\cPrior}{\mathFix{c}}
\newcommand{\like}{Q}
\newcommand{\sadd}[1]{\mathFix{#1^{*}}}
\newcommand{\saddR}{\sadd{R}}
\newcommand{\saddP}{\sadd{P}}
\newcommand{\est}{T}

\newcommand{\fun}[1][]{\ifthenelse{\isempty{#1}}{}{(#1)}}

Consider a random variable $X$.
In probabilistic modeling, our knowledge about
$X$ is typically represented by a probability distribution, which serves as a basis for making inferences.

However, in many situations, our information may be too limited to justify choosing a single distribution.
Instead, we may only be able to specify a set of plausible distributions.
This set can be viewed as a representation of \emph{partial knowledge}, or equivalently, a state of lack of
knowledge or \emph{ignorance}.

A well-known example of this idea is the exchangeability assumption used in de Finetti’s theorem.
A sequence of random variables \infSeqExp{X} is said to be exchangeable, if the ordering of the variables carries no
information.
That is, the joint probability distribution remains invariant under permutations of the variables.
This assumption imposes a structural constraint, effectively narrowing the set of admissible distributions.
Thus, exchangeability can be seen as a specific form of partial knowledge.

In this work, we adopt the perspective that such sets provide a meaningful representation of the state of partial
knowledge, and that this representation can likewise be used for inference.

Importantly, the act of choosing this set reflects some degree of knowledge.
If even that is uncertain, we may represent our uncertainty by a set of sets of distributions---capturing deeper
layers of ignorance.

In Bayesian theory, any type of knowledge must ultimately be encoded in a single probability measure.
Thus, to use our partial knowledge for performing inference in a Bayesian manner, we must find a principled way to
convert our set-valued representation into a single probability measure.

The probability measure that we choose, represents our beliefs about the environment.
It can be argued that this choice must depend solely on the environment.
As long as the environment is the same, we must not change this choice.
For example, if our loss function changes, or if the evidence that we get from the environment becomes more or
less informative, none of these should change our beliefs about the environment, and therefore our chosen
probability measure.
We will use this principle several times in this discussion and refer to it as the \emph{principle of objectivity}.

Subjectively, for a set of distributions, a probability measure that has the lowest divergence from the distributions
of the set can be considered the ``safest'' choice.
``Safe'' in the sense that it has the best possible performance in the worst case, and more importantly, when
convergence is possible for a set, if we use this probability measure, we are guaranteed to have a convergent
distribution.

\begin{definition}[Central Distribution]
    Let $X$ be a random variable and \setQ be a set of probability distributions over $X$.
    We say a probability distribution $C$ is a \emph{central distribution} of \setQ for $X$, if for any arbitrary
    distribution $P$, we have
    \[
        \supD[\cDist] \leq \supD[P] ,
    \]
    where $\measure{D}{\mathrm{KL}}{}$ denotes the Kullback–Leibler (KL) divergence.
\end{definition}

Note that with this definition the central distribution of \setQ is not essentially a member of \setQ.

\begin{definition}[Interior Distribution]
    Let $R, Q \in \setQ$ be such that for every $P \in \setQ$, we have
    \begin{equation*}
        \kl[R] \leq \kl[Q] \cdot
    \end{equation*}
    We refer to $R$ as an \emph{interior distribution} of \setQ.
\end{definition}

\begin{theorem}
    Let $R$ be an interior distribution of \setQ.
    If $C$ is a central distribution of $\setQ - \{R\}$, it is also a central distribution of \setQ, and vice versa.
\end{theorem}

\begin{theorem}
    Any distribution $R \in \setQ$ that can be expressed as a convex combination of elements of \setQ, is an
    interior distribution of \setQ.
\end{theorem}

\begin{theorem}
    Let $X$ and $T$ be two random variables, and let \setQ be a set of distributions over the pair $(X,T)$.
    Assume that for every \inSet{Q}, we have
    \[
        Q(X, T) = P(\XIT) Q(T),
    \]
    where $P$ is a fixed conditional distribution.
    If $C$ is a central distribution of \setQ for $T$, then $P(\XIT)C(T)$ is a central
    distribution of \setQ for $(X,T)$.
\end{theorem}

\begin{definition}[Central Divergence]
    For any distribution $P$, we define the \emph{central divergence} of $P$ from a set of distributions \setQ, with
    respect to a random variable $X$, as follows:
    \begin{equation*}
        \myDiv = \inf_{\cDist} \kl[{\cDist}],
    \end{equation*}
    where $\cDist$ denotes a central distribution of \setQ for $X$.
\end{definition}

When \setQ has a unique central distribution, the central divergence is simply defined as the KL divergence from
that unique central distribution.

\begin{definition}
    Let $X$ be a random variable, and let \setQQ be a set whose elements are sets of distributions over $X$.
    Let \cDist[\setQ] be a central distribution of $\setQ \in \setQQ$ for $X$, and $P$ be an arbitrary probability
    distribution.
    We say \ccDist is a central distribution of \setQQ, if we have:
    \begin{equation*}
        \sup_{\setQ \in \setQQ} \myDiv[\ccDist]  \leq \sup_{\setQ \in \setQQ} \myDiv \cdot
    \end{equation*}
\end{definition}

\begin{definition} [Saddle Point]
    Consider a function $f:\set{w} \times \set{z} \to \field{R}$.
    We refer to a pair $(w^*, z^*)$ as a \emph{saddle point} for $f$, if we have
    \[
        f(w^*, z) \leq f(w^*, z^*) \leq f(w, z^*),
    \]
    for all \inSet{w} and \inSet{z}.
\end{definition}

\begin{theorem}
    \label{th:saddle_is_central}
    Let $X$ be a discrete random variable, and \setQ be a finite set of distributions for $X$.
    Consider the following function:
    \begin{equation*}
        f(P, R)= \fExp,
    \end{equation*}
    where $Q_i \in \setQ$, and $P$, $R$ are probability distributions.

    If \saddle is a saddle point of $f$, then \saddP is a central distribution of \setQ for $X$.
\end{theorem}

\begin{proof}
    Since $X$ is a discrete random variable, \setQ has a central distribution for $X$.
    Let \cDist be a central distribution of \setQ, and \saddle be a saddle point for $f$.
    Clearly, the point $(\cDist, \saddR)$, is in the domain of $f$, and
    since \saddle is a saddle point for $f$, we have
    \begin{equation}
        \label{eq:th_proof_saddle_lq_dx}
        f\saddle \leq f(\cDist, \saddR) \cdot
    \end{equation}
    Because \saddR is a probability distribution, the definition of $f$ implies:
    \[
        f(\cDist, \saddR) \leq \supD[\cDist] \cdot
    \]
    On the other hand,
    \[
        f\saddle = \max_{R} f(\saddP, R) = \supD[\saddP] \cdot
    \]
    Thus, Inequality~\ref{eq:th_proof_saddle_lq_dx} implies
    \[
        \supD[\saddP] \leq \supD[\cDist] \cdot
    \]
    Consequently, $\saddP$ is also a central distribution of \setQ for $X$.
\end{proof}

\begin{lemma}
    \label{lm:saddle_conditions}
    Let $f$ be the function defined in theorem~\ref{th:saddle_is_central}.
    The point \saddle is a saddle point for $f$, if and only if there are constants $\lambda$, $\eta$ and a
    function $\mu$, such that for every $i$, the following conditions hold:
    \begin{gather}
        \label{eq:kkt_lambda}
        \klExp[Q_i][\saddP][x] - \mu(i) = \lambda, \\
        \label{eq:kkt_w_avg}
        \saddP(x) = \sum_i \saddR(i) Q_i(x) \\
        \label{eq:kkt_r_sums_one}
        \sum_i \saddR(i) = 1, \\
        \label{eq:saddle_lm_r_geq_0}
        \saddR(i) \geq 0, \\
        \nonumber
        \mu(i) \geq 0, \\
        \label{eq:saddle_lm_rmu_is_zero}
        \mu(i) \saddR(i) = 0 \cdot
    \end{gather}
\end{lemma}

\begin{proof}
    Similar to $f$, we define $\hat{f}$ as follows:
    \begin{equation*}
        \hat{f}(P, R) = \fExp,
    \end{equation*}
    where $P$, $R$ are arbitrary real functions instead of probability distributions.
    This enables us to formulate the saddle point of $f$ as the solution to a constrained optimization problem.

    Let \saddle be a simultaneous solution to the following two optimization problems:

    \begin{equation}
        \begin{aligned}
            &\text{minimize} & &\hat{f}(P,R) & &\text{w.r.t.} \ P \\
            &\text{subject to} & &\sum_{x} P(x) = 1, \\
            \\
            &\text{maximize} & &\hat{f}(P,R) & &\text{w.r.t.} \ R \\
            &\text{subject to} & &R(i) \geq 0, & &\text{for all} \ i  \\
            & & &\sum_{i} R(i) = 1 &
        \end{aligned}\label{eq:twin_optimization}
    \end{equation}
    Clearly, if $\saddle$ exists, it will be a saddle point for $f$.

    Both of these optimization problems are convex~\cite{0521833787} and differentiable.
    To establish this, consider that all the constraint functions are affine, and therefore are convex and
    differentiable.

    Additionally, $\hat{f}$ is differentiable function with respect to both $R$ and $P$.
    For any fixed $P$, $\hat{f}(P,\, \cdot \,)$ is an affine function and is concave.

    On the other hand, for a fixed solution to the second optimization problem denoted by $\hat{R}$, we can write
    \begin{equation}
        \label{eq:simplified_f}
        \hat{f}(P,\hat{R}) = c - \sum_{i,x} \hat{R}(i) Q_i(x) \ln P(x),
    \end{equation}
    where $c$ is not a function of $P$.
    Since $\hat{R}$ is a solution to the second optimization problem, it satisfies the constraints
    of the second problem and therefore is a probability distribution.
    Hence, for all $i$, $x$ we have $\hat{R}(i)Q_i(x) \geq 0$, and $\hat{f}(P,\, \cdot \,)$  is
    the negative of a linear combination of concave $\ln P(x)$ functions with nonnegative coefficients.
    This function is convex.

    Thus, both optimization problems are convex and differentiable, with affine constraints.
    The constraints simply restrict the feasible set to the set of probability distributions, and it is always possible
    to find feasible points that meet all the constraints.
    Therefore, Slater’s condition is satisfied, and the KKT conditions provide necessary
    and sufficient conditions for optimality~\cite{0521833787}.
    In other words, If a point $\saddle$ satisfies the KKT conditions for both problems, it will be a saddle
    point for $f$, and vice versa.

    The KKT conditions are as follows for every $x$, $i$ and constants $\lambda$, $\eta$:
    \begin{gather}
        \label{eq:kkt_proof_R_is_pxIe}
        \sum_i \saddR(i) Q_i(x) = \eta \saddP(x) \\
        \label{eq:kkt_proof_p_sums_one}
        \sum_{x} \saddP(x) = 1, \\
        \nonumber
        \klExp[Q_i][\saddP][x] - \mu(i) = \lambda, \\
        \label{eq:kkt_proof_r_sums_one}
        \sum_i \saddR(i) = 1, \\
        \nonumber
        \saddR(i) \geq 0, \\
        \nonumber
        \mu(i) \geq 0, \\
        \nonumber
        \mu(i) \saddR(i) = 0 \cdot
    \end{gather}

    By summing Equation~\ref{eq:kkt_proof_R_is_pxIe} over $x$ and using Equation~\ref{eq:kkt_proof_r_sums_one},
    we have
    \[
        \eta = 1 \cdot
    \]
    Thus, we get Equation~\ref{eq:kkt_w_avg} which also implicitly implies Equation~\ref{eq:kkt_proof_p_sums_one}.
\end{proof}

\begin{theorem} [Centrality]
    \label{th:lambda}
    Let $X$ be a discrete random variable, and let \setQ be a finite set of distributions over $X$.
    Suppose \cDist is a convex combination of elements of \setQ, with all weights strictly positive.
    If for every \inSet{Q} and some constant $\lambda$ we have:
    \begin{equation*}
        \kl[][\cDist] = \lambda,
    \end{equation*}
    then \cDist is a central distribution of \setQ for $X$.
\end{theorem}

\begin{proof}
    Lemma~\ref{lm:saddle_conditions} provides a set of sufficient conditions for a saddle point.
    Note that more restrictive conditions can also be imposed to derive an alternative set of sufficient conditions.
    We can replace Inequality~\ref{eq:saddle_lm_r_geq_0} with the more
    restrictive $\saddR(i) > 0$ condition.
    By Equation~\ref{eq:saddle_lm_rmu_is_zero}, that implies $\mu(i) = 0$,
    and Lemma~\ref{lm:saddle_conditions} conditions can be simplified to
    \begin{gather*}
        \klExp[Q][\saddP][x] = \lambda, \\
        \saddP(x) = \sum_i \saddR(i) Q(x),
    \end{gather*}
    which must hold for every \inSet{Q} and a positive probability distribution \saddR.

    Since \cDist satisfies these condition, Lemma~\ref{lm:saddle_conditions} and
    Theorem~\ref{th:saddle_is_central} imply that \cDist is a central distribution of \setQ for $X$.
\end{proof}

These theorems were proved under the assumption that \setQ is finite; however, they can be extended to cases where
\setQ can be indexed by a finite-dimensional real vector \param.
We can treat this vector as a random variable and interpret $Q_{\paramVal}$ as a conditional
distribution $Q(\xIt)$.
Then, the wight function, $\saddR(i)$, will essentially become a prior distribution for \paramVal.

\begin{theorem}
    Let $X$ be a discrete random variable and \setQ be a set of distributions over $X$, indexed by a
    finite-dimensional random vector \param.
    Let \cPrior be a positive prior distribution for \param such that for $\like(\xIt) \in \setQ$ we have
    \begin{equation*}
        \cPrior(\paramVal) \propto \exp \left\{  \sum_{x} \like(\xIt) \ln \cPrior(\tIx) \right\} \cdot
    \end{equation*}
    Then, $\cDist(x) = \mrg$ is a central distribution of \setQ for $X$.
    We also refer to \cPrior as a central prior for \param.
\end{theorem}

\begin{proof}
    for some constant $\eta$ we have
\end{proof}

Suppose we wish to model an exchangeable sequence of random variables.
Let $X_i$ denote the sequence of the first $i$ variables.
As discussed earlier, exchangeability can be interpreted as a form of partial knowledge, which allows us to
represent our uncertainty using a central distribution over $X_i$.

However, this central distribution generally depends on the sequence length $i$.
According to the principle of objectivity, our beliefs should not depend on the specifics of our experimental setup,
and in particular, should not vary with the length of the sequence.
It is important to note that a distribution defined over a longer sequence induces marginal distributions over its
shorter subsequences.

This observation suggests that, in order to satisfy the principle of objectivity, we may instead define a central
distribution over an infinite exchangeable sequence---thereby maintaining consistency with the
distributions of all finite initial segments.
This requires constructing a central distribution that is well-defined over $X_{\infty}$.

\begin{definition}
    Let \infSeq{X} be a sequence of random variables and \setQ be a set of probability measures that define
    the distribution of each $X_i$.
    Assume \cDist[i] is a central distribution for $X_i$.
    Let \cDist be a probability measure that defines a distribution over every $X_i$.
    \cDist is a central probability measure of \setQ for \infSeq{X}, if we have
    \begin{equation*}
        \lim_{i \to \infty} \kl[{\cDist[i]}][\cDist][X_i] = 0 \cdot
    \end{equation*}
\end{definition}

Intuitively, this definition states that \cDist is a central distribution for $X_{\infty}$.

\begin{theorem}
    Let \infSeq{\est} be a sequence of random variables, and let \setQ be a set of probability measures that define the
    distribution of each $\est_i$.
    Assume that \setQ can be indexed by a finite-dimensional vector \param.
    If \infSeq{\est} is a consistent sequence of estimators for \param, and there exists a sequence of
    constants \infSeq{\lambda} such that the sequence
    $\infSeq[]{\lambda_i -\entropy{Q}{\est_i \mid \paramVal}}$ converges
    for every value of \paramVal, then
    \begin{equation*}
        \cPrior(\paramVal) \propto \lim_{i \to \infty} \exp
        \left\{ \lambda_i -\entropy{Q}{\est_i \mid \paramVal} \right\},
    \end{equation*}
    is a central prior of \setQ for \param.
\end{theorem}
